\documentclass[10pt,a4paper]{amsart}
\usepackage[margin=19mm]{geometry}
\usepackage{amsmath,amssymb,mathtools}
\usepackage[scr=rsfs,cal=euler]{mathalfa}
\usepackage{xfrac}
\usepackage[unicode,hidelinks,bookmarks=false]{hyperref}
\newcommand{\ie}{i.e.\ }
\newcommand{\cf}{cf.\ }
\newcommand{\resp}{respectively\ }
\newcommand{\Subsection}{\subsection}
\providecommand{\nobreakdash}{\nobreak\hbox{-}}
\setlength{\emergencystretch}{2em}
\multlinegap0pt
\usepackage{bm}
\usepackage{bbm}
\usepackage{dsfont}
\usepackage{xspace}
\usepackage{cancel}
\usepackage{mathtools}
\usepackage{amsthm}
\makeatletter
\@ifundefined{theorem}{\newtheorem{theorem}{Theorem}}{}
\@ifundefined{lemma}{\newtheorem{lemma}[theorem]{Lemma}}{}
\makeatother
\DeclareMathOperator{\cl}{cl}
\def\e{\varepsilon}
\def\ov{\overline}
\def\sepa{\; : \;}
\newcommand{\RR}{\mathbb{R}}
\newcommand{\aco}[1]{\left\{#1\right\}}
\newcommand{\tinm}[1]{\text{\;#1\;}}
\title[Calibration for minimal surfaces with free boundary and Chee]{Calibration for minimal surfaces with free boundary and Cheeger-type problems}
\author{Guy Bouchitt\'e, Minh Phan}
\date{Source: arXiv:2511.03322v1 (CC BY 4.0)}
\begin{document}
\maketitle

\noindent\emph{Source and licence.} Calibration for minimal surfaces with free boundary and Cheeger-type problems. Version arXiv:2511.03322v1 (CC BY 4.0), distributed under the Creative Commons Attribution 4.0 International licence, as recorded in the arXiv licence metadata for this exact version and shown on the abstract page https://arxiv.org/abs/2511.03322v1. Full rights notice: licenses/arxiv-2511.03322-CC-BY-4.0.txt.\par\smallskip

\noindent\textit{Editorial scope.} Counted unit: the Lemma whose source label is \texttt{access} in the article (source0.tex lines 3405--3416, source label \texttt{access}), delivered together with the article's own complete original proof of it (source0.tex lines 3418--3471, one \texttt{proof} environment of 54 lines). No mathematics is written, repaired or re-derived: statement and proof are the article's own text line for line. Required context retained uncounted: none. Every definition, display and label of those lines is retained unchanged. Omitted from this package and not redistributed: 1--3404, 3417--3417, 3472--4100. Nothing inside a retained excerpt is omitted or altered: every retained body is delivered exactly as the article's lines are written, so no omission receipt of any kind accompanies this package. Citations: the retained excerpts carry no citation command at all, so no bibliography entry of the article is transcribed and no reference is redistributed. Reference adaptation: none was needed and none was made; every reference inside the delivered unit resolves inside it, so no unresolved reference or citation is produced. Printed slips kept literal: no printed word, symbol, hypothesis or number of the counted statement or of its proof was altered. Layout and class: the article's own class is \texttt{amsart} with packages including fontenc, inputenc, amsmath, amsfonts, amssymb, fullpage, latexsym, color, lmodern, enumitem, graphicx, tikz, dsfont, microtype; the wrapper is the lane's pinned \texttt{amsart} editorial wrapper at 10pt on A4, which restates the theorem-like environments the retained text needs and adds the article's own macro definitions that the retained text needs (\texttt{\textbackslash RR}, \texttt{\textbackslash aco}, \texttt{\textbackslash cl}, \texttt{\textbackslash e}, \texttt{\textbackslash ov}, \texttt{\textbackslash sepa}, \texttt{\textbackslash tinm}), with extra wrapper packages (bm, bbm, dsfont, xspace, cancel, mathtools). The delivered numbering is the unit's own. Assets: no retained span contains a figure environment, a graphics-inclusion command, a PDF-page inclusion, a table or a TikZ picture; no image file is copied into this package, the package declares no asset dependency and no image file is redistributed with it. Licence: version arXiv:2511.03322v1 of this article is distributed under the Creative Commons Attribution 4.0 International licence, as recorded in the arXiv licence metadata for this exact version and shown on the abstract page https://arxiv.org/abs/2511.03322v1. A full rights notice is delivered at licenses/arxiv-2511.03322-CC-BY-4.0.txt. Characters: every retained excerpt is pure ASCII, so the delivered engine, XeLaTeX with its default Latin Modern text font, reports no missing character.
\begin{lemma}\label{access}
Let $D$ be a convex  domain  of $\RR^2$. For every $x\in \partial D$, we define
\begin{align*}
k_{\partial D}(x):= \frac{1+ \nu_D^+(x)\cdot\nu_D^-(x)}{2}.
\end{align*}
Then, we have
\begin{itemize}
\item[(i)] for every $x\in\partial D$, $0<k_{\partial D}(x) \le 1$;
\item[(ii)] $k_{\partial D}(x) = 1$ for every $x\in \partial_r D$;
\item[(iii)] $\forall\e \in(0,1)$, the set  $\aco{x \sepa k_{\partial D}(x) <\e }$ is finite.
\end{itemize}
\end{lemma}

\begin{proof}
We recall that the normal cone of $D$ at $x$ is given by
\begin{align*}
N_{D}(x):=\aco{ a \nu_D^+(x) + b \nu_D^-(x) \sepa a,b \in \RR_+ }.
\end{align*}
For every $x\in \partial D$, we denote by $\varphi(x)$ the angle
\begin{align*}
\varphi(x):= \frac{1}{2}\angle(\nu_D^-(x),\nu_D^+(x)),
\end{align*}
and by $T_D(x)$ the tangent cone of $D$ at $x$
\begin{align*}
T_D(x):= & \cl \aco{ s(y-x) \sepa y\in D,\; s\ge 0 }\\
=& \aco{a T_D^+(x) + b T_D^-(x) \sepa a,b \in \RR_+},
\end{align*}
where $T_D^\pm(x)$ are the left and right tangent unit vectors of $D$ at $x$. We denote by $\psi(x)$ the angle 
\begin{align*}
\psi(x):= \frac{1}{2}\angle (T_D^-(x),T_D^+(x)).
\end{align*} 
Since tangent and normal cones are polar each other, it holds $\varphi(x)+\psi(x) = \frac{\pi}{2}$ for every $x\in \partial D$.


\noindent
\emph{Proof of (i):} clearly, as $D$ is convex, we have 
$ 0\le \varphi(x)< \frac{\pi}{2}$ and  $0< \psi(x)\le \frac{\pi}{2}$.
Hence, for all $x\in\partial D$, it holds $0<k_{\partial D}(x) \le 1$ since
\begin{align*}
k_{\partial D}(x)=\frac{1+\cos 2\varphi(x) }{2} =\cos^2\varphi(x) = \sin^2 \psi(x).
\end{align*}

\noindent
\emph{Proof of (ii):}  As $D$ is convex, its boundary $\partial D$ admits at most countably many singular points. On the regular part $\partial_r D=\partial D\setminus \partial_s D$, we get $\nu_D^-(x)=\nu_D^+(x)$, i.e. $\varphi(x)=0$. Then, $k_{\partial D}(x)=1$.

\medskip
\noindent
\emph{Proof of (iii):} Let $\e$ such that $0<\e<1$ and denote
\begin{align*}
E_\e:=\{ x\in\partial D \sepa k_{\partial D}(x) <\e \}\quad,\quad N_\e:=\#(E_\e) .
\end{align*}
  The function $h(\varphi)$ defined by
\begin{align*}
h(\varphi):=\frac{1+\cos 2\varphi}{2}, \quad \tinm{for} \varphi\in [0,\frac{\pi}{2}).
\end{align*}
 is non increasing in $[0,\frac{\pi}{2})$, while  $h(\varphi(x))=k_{\partial D}(x)$ for every $x\in \partial D$.
There always exists $\ \varphi_e\in (0,\pi/2)$ such that $h(\varphi_\e)=\e$. We set
\begin{align*}
E_\e:=\{ x\in\partial D \sepa k_{\partial D}(x) <\e \}.
\end{align*}
By the monotonicity of $h$, we deduce that  $\varphi(x) > \ov \varphi$ for every $x\in E_\e$.
It follows that
\begin{align*}
N_\e\, \ov\varphi  < \sum_{x\in E}\varphi(x) \le \sum_{x\in\partial D}\varphi(x) \le 2\pi ,
\end{align*}
whence $N_\e \le \frac{2\pi}{ \ov\varphi} <+\infty$.
\end{proof}

\end{document}
